\subsection*{Intended Learning Outcomes - Lecture 08}

The student should be able to:
\begin{itemize}
\item write down the 3 do-calculus rules and reason with them;
\item prove the 3 rules of do-calculus using the global Markov property;
\item apply the 3 rules of do-calculus in an iterative manner on (small, not too complicated) graphs for the identification of causal effects (i.e. reducing the number $\doit$-terms);
\item state and apply the back-door covariate adjustment criterion and formula;
\item know and understand the definition and consequences of unique solvability for SCMs;
\item know and understand the definition and consequences of an SCM being simple;
\item be able to check whether a given SCM is uniquely solvable/simple;
\item be able to reason about counterfactuals by using the twin SCM;
\item understand the different equivalence relations for SCMs;
\item know the basic commutation and compatibility properties of basic operations on SCMs (twinning, intervening, preservation of unique solvability / simplicity / equivalence)
\end{itemize}
